\documentclass[12pt,a4paper]{article}
\usepackage[utf8]{inputenc}
\usepackage[T2A]{fontenc}
\usepackage[russian]{babel}
\usepackage{amsmath,amssymb,amsthm}
\usepackage{booktabs,array}
\usepackage{microtype}
\usepackage{geometry}
\geometry{margin=2.3cm}
\setlength{\emergencystretch}{2em}

\newtheorem{proposition}{Предложение}
\newtheorem{nogocriterion}{Критерий провала}

\title{Меню состояний с минимальным неабелевым fiber\\
Проверка \(SU(5)\)-расширения}
\author{}
\date{4 августа 2026 г.}

\begin{document}
\maketitle

\section{Постановка gate}

Первый state-menu audit показал, что
\[
\pi_1(\mathbb{RP}^3\times S^1)=\mathbb Z_2\times\mathbb Z
\]
даёт две фазовые окружности и метки \((\varepsilon,n)\), но остаётся полностью абелевым. Поэтому следующий заранее объявленный вопрос имеет вид:
\begin{quote}
Существует ли один минимальный простой неабелев fiber над меню состояний, который одной \(\mathbb Z_2\)-голономией порождает одновременно color triplet и weak doublet, не вводя независимые \(SU(3)\)- и \(SU(2)\)-сектора?
\end{quote}

\section{Почему возникает \(SU(5)\)}

Ранг стандартной калибровочной группы равен
\[
\operatorname{rank}(SU(3)\times SU(2)\times U(1))=4.
\]
Следовательно, простой compact fiber должен иметь ранг не меньше четырёх. Среди простых групп ранга четыре минимальную размерность имеет
\[
SU(5),
\qquad
\dim SU(5)=24.
\]
Для сравнения, \(SO(8)\) имеет размерность \(28\), \(SO(9)\) и \(Sp(8)\) --- \(36\), а \(F_4\) --- \(52\).

\begin{proposition}[Минимальность fiber в rank--dimension gate]
В классе compact connected simple Lie groups, способных содержать rank-four subgroup стандартной модели, \(SU(5)\) является минимальным по размерности Lie algebra кандидатом.
\end{proposition}

\section{Торсионная голономия}

Нетривиальный генератор \(\mathbb Z_2\) можно отобразить в
\[
P=\operatorname{diag}(1,1,1,-1,-1)\in SU(5).
\]
Действительно,
\[
P^2=I,
\qquad
\det P=1.
\]
Коммутант этой голономии сохраняет трёхмерный и двумерный blocks:
\[
C_{SU(5)}(P)=S(U(3)\times U(2))
\cong\frac{SU(3)\times SU(2)\times U(1)}{\mathbb Z_6}.
\]
Его размерность равна
\[
8+3+1=12,
\]
а оставшиеся двенадцать генераторов \(SU(5)\) становятся broken directions.

Таким образом, прежняя дискретная ветвь меню не просто допускает неабелев fiber, а имеет ровно тот order-two элемент, который разбивает минимальный fiber до калибровочной группы стандартной модели.

\section{Одно поколение}

Возьмём стандартный traceless hypercharge generator
\[
Y=\operatorname{diag}\left(-\frac13,-\frac13,-\frac13,
\frac12,\frac12\right).
\]
Минимальный chiral anomaly-free package равен
\[
\mathbf{10}\oplus\overline{\mathbf5}.
\]
Он раскладывается как
\begin{align}
\mathbf{10}
&\longrightarrow
(\mathbf3,\mathbf2)_{1/6}
\oplus(\overline{\mathbf3},\mathbf1)_{-2/3}
\oplus(\mathbf1,\mathbf1)_1,\\
\overline{\mathbf5}
&\longrightarrow
(\overline{\mathbf3},\mathbf1)_{1/3}
\oplus(\mathbf1,\mathbf2)_{-1/2}.
\end{align}
Это ровно набор \(Q,u^c,e^c,d^c,L\) одного поколения без правого нейтрино.

Кубическая \(SU(5)\)-аномалия сокращается:
\[
\mathcal A(\mathbf{10})+\mathcal A(\overline{\mathbf5})
=1-1=0.
\]
После branching также сокращаются mixed \(SU(3)^2U(1)\), \(SU(2)^2U(1)\), gravitational--\(U(1)\) и \(U(1)^3\) anomalies.

\section{Неожиданный parity result}

Голономия \(P\) действует со знаком, равным произведению знаков fundamental indices. Поэтому получаем:
\begin{center}
\begin{tabular}{>{\raggedright\arraybackslash}p{0.42\textwidth}c}
\toprule
Состояния & \(\mathbb Z_2\)-parity \\
\midrule
Weak doublets \(Q\) и \(L\) & \(-1\) \\
Weak singlets \(u^c,d^c,e^c\) & \(+1\) \\
\bottomrule
\end{tabular}
\end{center}

\begin{proposition}[Единый weak-doublet selector]
Одна и та же торсионная голономия, разбивающая \(SU(5)\) до стандартной калибровочной группы, делает все weak doublets нечётными, а все weak singlets чётными. Независимый sector-by-sector parity assignment не требуется.
\end{proposition}

Это первый действительно положительный cross-sector результат новой menu-гипотезы.

\section{Соблазн трёх поколений}

На \(K\) существуют четыре spin-структуры. Поэтому возникает очевидный соблазн:
\[
4-1=3,
\]
то есть выбрать один vacuum/reference sector и отождествить три остальных с поколениями.

Однако spin-структуры образуют affine torsor над \(\mathbb Z_2^2\). Без дополнительного объекта в torsor нет канонического нуля и, следовательно, нет канонически выделенной структуры, которую следует исключить.

\begin{nogocriterion}[Провал автоматического generation count]
Равенство \(4-1=3\) является перспективной подсказкой, но не выводом трёх поколений. Требуется независимый геометрический инвариант, defect или boundary condition, единственным образом выделяющий reference spin-sector до сравнения с наблюдаемым числом поколений.
\end{nogocriterion}

\section{Вердикт и следующий gate}

Второй menu-gate проходит существенно лучше первого. Один минимальный \(SU(5)\)-fiber и одна уже существующая \(\mathbb Z_2\)-голономия дают:
\begin{itemize}
    \item стандартную калибровочную группу как centralizer;
    \item color triplets и weak doublets из одного fiber;
    \item anomaly-free chiral package одного поколения;
    \item единый parity rule, отделяющий doublets от singlets.
\end{itemize}

При этом не выведены выбор \(\mathbf{10}\oplus\overline{\mathbf5}\), три поколения, Yukawa matrices и массы. Следующий gate должен искать invariant четырёх spin-структур, который без знания числа поколений выделяет один reference sector и оставляет три эквивалентных nonreference sectors.

\end{document}